% Λύση του 55ου προβλήματος της ενότητας «Άλυτα Προβλήματα».
\subsection*{Λύση Προβλήματος 55 --- Λογιστικό μοντέλο πληθυσμού}

\begin{alist}
  \item Θέτουμε
        \[
          a=\frac{K-P_0}{P_0}.
        \]
        Επειδή \(0<P_0<K\), έχουμε \(a>0\). Η συνάρτηση γράφεται
        \[
          P(t)=\frac{K}{1+ae^{-rt}}.
        \]
        Παραγωγίζουμε:
        \[
          P'(t)
          =
          K\left(1+ae^{-rt}\right)^{-1}'.
        \]
        Με τον κανόνα σύνθετης συνάρτησης,
        \[
          P'(t)
          =
          -K\left(1+ae^{-rt}\right)^{-2}
          \cdot
          \left(1+ae^{-rt}\right)'.
        \]
        Επειδή
        \[
          \left(1+ae^{-rt}\right)'=-ar e^{-rt},
        \]
        παίρνουμε
        \[
          P'(t)
          =
          \frac{Kare^{-rt}}{\left(1+ae^{-rt}\right)^2}.
        \]
        Άρα
        \[
          {
          P'(t)=\frac{Kare^{-rt}}{\left(1+ae^{-rt}\right)^2}
          }.
        \]

  \item Θα δείξουμε ότι
        \[
          P'(t)=rP(t)\left(1-\frac{P(t)}{K}\right).
        \]
        Από τον τύπο του \(P(t)\),
        \[
          \frac{P(t)}{K}
          =
          \frac{1}{1+ae^{-rt}}.
        \]
        Άρα
        \[
          1-\frac{P(t)}{K}
          =
          1-\frac{1}{1+ae^{-rt}}
          =
          \frac{ae^{-rt}}{1+ae^{-rt}}.
        \]
        Επομένως
        \[
          rP(t)\left(1-\frac{P(t)}{K}\right)
          =
          r\cdot
          \frac{K}{1+ae^{-rt}}
          \cdot
          \frac{ae^{-rt}}{1+ae^{-rt}}.
        \]
        Άρα
        \[
          rP(t)\left(1-\frac{P(t)}{K}\right)
          =
          \frac{Kare^{-rt}}{\left(1+ae^{-rt}\right)^2}.
        \]
        Αυτό είναι ακριβώς το \(P'(t)\). Συνεπώς
        \[
          {
          P'(t)=rP(t)\left(1-\frac{P(t)}{K}\right)
          }.
        \]

  \item Ο παράγοντας
        \[
          1-\frac{P}{K}
        \]
        εκφράζει τον περιορισμό που επιβάλλει το περιβάλλον στην
        ανάπτυξη του πληθυσμού.

        Όταν \(P\) είναι πολύ μικρό σε σχέση με το \(K\), τότε
        \[
          \frac{P}{K}\approx0
          \quad\Rightarrow\quad
          1-\frac{P}{K}\approx1.
        \]
        Άρα ο πληθυσμός αυξάνεται σχεδόν εκθετικά.

        Όταν \(P\) πλησιάζει το \(K\), τότε
        \[
          \frac{P}{K}\approx1
          \quad\Rightarrow\quad
          1-\frac{P}{K}\approx0.
        \]
        Άρα ο ρυθμός ανάπτυξης μικραίνει, επειδή οι πόροι του
        περιβάλλοντος δεν επαρκούν για απεριόριστη αύξηση.

  \item Ο ρυθμός ανάπτυξης, ως συνάρτηση του πληθυσμού \(P\), είναι
        \[
          R(P)=rP\left(1-\frac{P}{K}\right).
        \]
        Δηλαδή
        \[
          R(P)=rP-\frac{r}{K}P^2.
        \]
        Πρόκειται για παραβολή ως προς \(P\), με αρνητικό συντελεστή στο
        \(P^2\). Άρα έχει μέγιστο στην κορυφή της.

        Η κορυφή βρίσκεται στο
        \[
          P=-\frac{\beta}{2a_{\pi}},
        \]
        όπου εδώ
        \[
          a_{\pi}=-\frac{r}{K},
          \qquad
          \beta=r.
        \]
        Άρα
        \[
          P
          =
          -\frac{r}{2\left(-\frac{r}{K}\right)}
          =
          \frac{K}{2}.
        \]
        Επομένως ο ρυθμός ανάπτυξης είναι μέγιστος όταν
        \[
          {P(t)=\frac{K}{2}}.
        \]
        Δηλαδή η αποικία αυξάνεται ταχύτερα όταν έχει φτάσει στο μισό
        της χωρητικότητας του περιβάλλοντος.
\end{alist}
